% PAGES: 306-308

\begin{theorem}
(i) Matrix multiplication by a permutation matrix to the left results in a
corresponding permutation of the rows of the matrix.

In other words, if $A = \mathrm{row}(r_j)_{j=1:n}$ is a square matrix of size $n$ and
if $P = \mathrm{row}\,\!\left(e_{\tau_r(j)}^t\right)_{j=1:n}$ is a permutation matrix,
then
\begin{equation}
PA \;=\; \mathrm{row}\left(r_{\tau_r(j)}\right)_{j=1:n}.
\end{equation}

(ii) Matrix multiplication by a permutation matrix to the right results in a
corresponding permutation of the columns of the matrix.

In other words, if $A = \mathrm{col}(a_k)_{k=1:n}$ is a square matrix of size $n$ and
if $P = \mathrm{col}\,\!\left(e_{\tau_c(k)}\right)_{k=1:n}$ is a permutation matrix,
then
\begin{equation}
AP \;=\; \mathrm{col}\left(a_{\tau_c(k)}\right)_{k=1:n}.
\end{equation}
\end{theorem}

\begin{proof}
(i) Recall from (1.27) that $e_j^tA = r_j$ for all $j = 1:n$. Thus,
\[
e_{\tau_r(j)}^tA \;=\; r_{\tau_r(j)}, \quad \forall\, j = 1:n,
\]
and therefore, from the matrix--matrix multiplication formula (1.12), we find that
\[
PA \;=\; \mathrm{row}\left(e_{\tau_r(j)}^tA\right)_{j=1:n} \;=\; \mathrm{row}\left(r_{\tau_r(j)}\right)_{j=1:n}.
\]

(ii) Recall from (1.26) that $Ae_k = a_k$ for all $k = 1:n$. Thus,
\[
Ae_{\tau_c(k)} \;=\; a_{\tau_c(k)}, \quad \forall\, j = 1:n,
\]
and therefore, from the matrix--matrix multiplication formula (1.11), we find that
\[
AP \;=\; \mathrm{col}\left(Ae_{\tau_c(k)}\right)_{k=1:n} \;=\; \mathrm{col}\left(a_{\tau_c(k)}\right)_{k=1:n}.
\qedhere
\]
\end{proof}

\begin{examplex}
Let
\[
A \;=\; \begin{pmatrix}
-2 & 3 & -1 & 0 & 15 \\
1 & 0 & -10 & -2 & 4 \\
5 & -2 & 1 & 0 & 1 \\
-5 & -1 & 2 & -6 & -3 \\
0 & -1 & -3 & 11 & -9
\end{pmatrix}
\]
and let $P$ be the following permutation matrix:
\[
P \;=\; \begin{pmatrix}
0 & 0 & 1 & 0 & 0 \\
1 & 0 & 0 & 0 & 0 \\
0 & 0 & 0 & 0 & 1 \\
0 & 1 & 0 & 0 & 0 \\
0 & 0 & 0 & 1 & 0
\end{pmatrix}.
\]

Then,
\[
PA \;=\; \begin{pmatrix} e_3^t \\ e_1^t \\ e_5^t \\ e_2^t \\ e_4^t \end{pmatrix} A
\;=\; \begin{pmatrix}
5 & -2 & 1 & 0 & 1 \\
-2 & 3 & -1 & 0 & 15 \\
0 & -1 & -3 & 11 & -9 \\
1 & 0 & -10 & -2 & 4 \\
-5 & -1 & 2 & -6 & -3
\end{pmatrix}
\]
\[
AP \;=\; (e_2 \mid e_4 \mid e_1 \mid e_5 \mid e_3)\, A \;=\; \begin{pmatrix}
3 & 0 & -2 & 15 & -1 \\
0 & -2 & 1 & 4 & -10 \\
-2 & 0 & 5 & 1 & 1 \\
-1 & -6 & -5 & -3 & 2 \\
-1 & 11 & 0 & -9 & -3
\end{pmatrix}
\]
\end{examplex}

\begin{lemma}
The product of two permutation matrices is a permutation matrix.
\end{lemma}

\begin{lemma}
(i) The determinant of any permutation matrix is either $1$ or $-1$. Thus, any
permutation matrix is nonsingular.

(ii) Any permutation matrix is orthogonal, i.e., the inverse of any permutation
matrix is its transpose. In other words, if $P$ is a permutation matrix, then
$P^{-1} = P^t$.
\end{lemma}

\subsection{Orthogonality}

\begin{definition}
Two vectors of the same size are orthogonal if and only if their inner product is
equal to $0$.

In other words, two vectors $u$ and $v$ of the same size are orthogonal if and only if
\[
(u,v) \;=\; v^tu \;=\; 0.
\]
\end{definition}

\begin{definition}
A square matrix is orthogonal\footnote{Throughout this book, we require, by
definition, that any column of an orthogonal matrix has norm equal to $1$. An
orthogonal matrix with the properties from Definition 10.3 is also called an
orthonormal matrix.} if and only if any two different columns of the matrix are
orthogonal and the norm of every column is equal to $1$.\footnote{Note that any two
rows of an orthogonal matrix are also orthogonal, and the norm of every row of an
orthogonal matrix is equal to $1$; cf. Lemma 10.5.}

In other words, the $n \times n$ matrix $Q = \mathrm{col}(q_k)_{k=1:n}$ is orthogonal
if and only if
\begin{align}
(q_k,q_j) \;=\; q_j^tq_k \;=\; 0, &\qquad \forall\, 1 \le j \ne k \le n; \\
||q_i||^2 \;=\; q_i^tq_i \;=\; 1, &\qquad \forall\, i = 1:n.
\end{align}
\end{definition}

\begin{theorem}
A square matrix is orthogonal if and only if its transpose matrix is also its inverse
matrix.

In other words, the square matrix $Q$ is orthogonal if and only if
\begin{equation}
Q^tQ \;=\; QQ^t \;=\; I,
\end{equation}
i.e., if and only if
\begin{equation}
Q^t \;=\; Q^{-1}.
\end{equation}
\end{theorem}

\begin{proof}
Let $Q = \mathrm{col}(q_k)_{k=1:n}$ be a square matrix of size $n$. Then,
\[
Q^t \;=\; \mathrm{row}\left(q_k^t\right)_{k=1:n} \;=\; \mathrm{row}\left(q_j^t\right)_{j=1:n}.
\]

Recall from (1.13) that the $(j,k)$ entry of the matrix $Q^tQ$ is given by
\begin{equation}
(Q^tQ)(j,k) \;=\; q_j^tq_k, \quad \forall\, 1 \le j,k \le n.
\end{equation}

Since the identity matrix $I$ has all entries equal to $0$ except for the entries on
its main diagonal which are equal to $1$, and using (10.16), (10.12), and (10.13), we
obtain that
\begin{align}
Q^tQ = I \quad &\iff \quad
\begin{cases}
(Q^tQ)(j,k) = 0, & \forall\, 1 \le j \ne k \le n \\
(Q^tQ)(i,i) = 1, & \forall\, i = 1:n
\end{cases} \notag \\
&\iff \quad
\begin{cases}
q_j^tq_k = 0, & \forall\, 1 \le j \ne k \le n \\
q_i^tq_i = 1, & \forall\, i = 1:n
\end{cases} \notag \\
&\iff \quad
\begin{cases}
(q_k,q_j) = q_j^tq_k = 0, & \forall\, 1 \le j \ne k \le n \\
||q_i||^2 = q_i^tq_i = 1, & \forall\, i = 1:n
\end{cases}
\end{align}

From (10.17) and Definition 10.3, we conclude that $Q^tQ = I$ if and only if the
matrix $Q$ is orthogonal.

Then, from Lemma 1.6, it follows that $Q^t = Q^{-1}$ and $QQ^t = I$.
\end{proof}

\begin{lemma}
A square matrix is orthogonal if and only if any two different rows of the matrix are
orthogonal and the norm of every row is equal to $1$.
\end{lemma}

\begin{proof}
Let $Q = \mathrm{row}(r_j)_{j=1:n}$ be an $n \times n$ orthogonal matrix. Then,
\[
Q^t \;=\; \mathrm{col}\left(r_j^t\right)_{j=1:n} \;=\; \mathrm{col}\left(r_k^t\right)_{k=1:n}.
\]

From Theorem 10.2, it follows that $Q$ is an orthogonal matrix if and only if
$QQ^t = I$. Note that the $(j,k)$ entry of $QQ^t$ is given by
\[
(QQ^t)(j,k) \;=\; r_jr_k^t, \quad \forall\, 1 \le j,k \le n;
\]
see (1.13). Then,
% No equation number printed for this chain in the source (contrast with the
% analogous chain in Theorem 10.2's proof, which does carry (10.17)) -- using
% align* to match, per fidelity rule (never normalize the book's own
% inconsistent numbering).
\begin{align*}
QQ^t = I \quad &\iff \quad
\begin{cases}
(QQ^t)(j,k) = 0, & \forall\, 1 \le j \ne k \le n; \\
(QQ^t)(i,i) = 1, & \forall\, i = 1:n.
\end{cases} \\
&\iff \quad
\begin{cases}
r_j^tr_k = 0, & \forall\, 1 \le j \ne k \le n; \\
r_j^tr_j = 1, & \forall\, i = 1:n.
\end{cases} \\
&\iff \quad
\begin{cases}
(r_k,r_j) = r_j^tr_k = 0, & \forall\, 1 \le j \ne k \le n; \\
||r_i||^2 = r_i^tr_i = 1, & \forall\, i = 1:n.
\end{cases}
\end{align*}

We conclude that the matrix $Q$ is orthogonal if and only if any two rows of $Q$ are
orthogonal and the norm of every row of $Q$ is equal to $1$.
\end{proof}

\begin{lemma}
The product of two orthogonal square matrices of the same size is an orthogonal
matrix.
\end{lemma}

\begin{proof}
Let $Q = Q_1Q_2$, where $Q_1$ and $Q_2$ are orthogonal matrices of the same size.
Since $Q_1^tQ_1 = I$ and $Q_2^tQ_2 = I$, we obtain that
\begin{align*}
Q^tQ \;&=\; (Q_1Q_2)^tQ_1Q_2 \;=\; Q_2^tQ_1^tQ_1Q_2 \\
&=\; Q_2^t(Q_1^tQ_1)Q_2 \;=\; Q_2^tQ_2 \\
&=\; I,
\end{align*}
and, from Theorem 10.2, we conclude that the matrix $Q$ is orthogonal.
\end{proof}
